Subbab~\ref{sec:pembahasan_terintegrasi} telah membahas secara terintegrasi
mengapa pola-pola pada hasil eksperimen muncul dan bagaimana pola tersebut
berkaitan dengan rancangan metode usulan. Subbab ini menutup BAB IV dengan
merangkum temuan utama dari delapan skenario, sepuluh \textit{trial} per
skenario, dan seluruh analisis deskriptif, statistik, serta kualitatif yang
telah dilakukan, sebagai jembatan menuju kesimpulan akhir pada BAB V.

\subsection{Ringkasan Temuan Utama}

Pada aspek keberhasilan navigasi, metode usulan mempertahankan tingkat
keberhasilan penuh pada seluruh skenario yang diujikan padanya, termasuk
pada kondisi \textit{obstacle} dinamis yang menjadi satu-satunya kondisi
\textit{baseline} mengalami kegagalan. Hal ini tidak berarti \textit{baseline}
gagal secara umum, karena \textit{baseline} tetap berhasil pada sebagian
besar skenarionya, namun metode usulan terbukti menutup celah kegagalan yang
muncul pada kondisi paling menantang.

Pada aspek keselamatan, minimum clearance meningkat secara konsisten, dan
peningkatan ini tidak merata, melainkan membesar seiring kompleksitas
skenario. Pada lingkungan sederhana tanpa \textit{obstacle}, perubahan tidak
relevan untuk dievaluasi; pada lingkungan dengan \textit{obstacle} statis
tunggal, peningkatan tergolong kecil; sementara pada kondisi yang melibatkan
\textit{blind-spot} dan \textit{obstacle} dinamis, peningkatan margin
keselamatan paling nyata. Pola ini menunjukkan bahwa kontribusi metode
usulan paling terasa justru pada kondisi yang memang menjadi sasaran
rancangannya, bukan pada kondisi navigasi yang sudah sederhana sejak awal.

Pada aspek efisiensi, peningkatan keselamatan tersebut tidak diperoleh
tanpa konsekuensi. Waktu tempuh meningkat pada seluruh pasangan skenario,
dengan besaran yang sebagian besar berbanding lurus dengan besarnya
peningkatan keselamatan, meskipun tidak selalu proporsional secara seimbang.
Panjang lintasan, sebaliknya, relatif tidak banyak berubah, sehingga
\textit{trade-off} yang terjadi lebih tepat dipahami sebagai konsekuensi
pada kecepatan dan durasi gerak, bukan pada rute geometris yang dipilih
robot. \textit{Trade-off} ini sebaiknya dibaca sebagai pertimbangan yang
melekat pada pendekatan yang lebih konservatif, bukan sebagai kelemahan
mutlak metode usulan.

Pada aspek kualitas gerak, tidak ditemukan indikasi bahwa peningkatan
keselamatan dicapai melalui perilaku navigasi yang ekstrem. Pemeriksaan
visual terhadap trajektori tidak menemukan pola zig-zag tajam pada
konfigurasi manapun, secara angular velocity, kecepatan angular metode
usulan bahkan lebih rendah pada sebagian skenario, meski pada skenario
lain justru sedikit lebih tinggi atau tidak berbeda secara berarti. Kualitas
gerak metode usulan dengan demikian tetap berada pada tingkat yang dapat
diterima, tanpa dapat diklaim secara seragam "lebih halus" pada seluruh
kondisi.

Secara statistik, hampir seluruh perbedaan yang diuji pada penelitian ini
signifikan, dan effect size yang menyertainya secara umum tergolong besar,
menunjukkan bahwa perbedaan-perbedaan tersebut bukan sekadar artefak
ukuran sampel. Meski demikian, signifikansi statistik pada penelitian ini
beberapa kali tidak berbanding lurus dengan besarnya selisih absolut yang
relevan secara praktis, sehingga interpretasi setiap hasil tetap dilakukan
dengan mempertimbangkan konteks aplikasi yang dituju, bukan semata dari
nilai $p$ atau $r$.

\subsection{Menjawab Tujuan Penelitian}

BAB I pada naskah ini belum memuat rumusan tujuan penelitian secara
definitif. Sebagai acuan kerja, tujuan yang secara konsisten mendasari
rancangan metode pada BAB III dan pengujian pada BAB IV ini adalah
mengevaluasi apakah integrasi \textit{virtual blind-spot scan}, fusi
\textit{traversability} LiDAR--termal, dan adaptasi fuzzy Mamdani dapat
meningkatkan keselamatan navigasi robot pada kondisi \textit{blind-spot}
dan lingkungan dinamis, tanpa mengorbankan performa secara berlebihan.
Berdasarkan acuan tersebut, hasil penelitian ini mendukung bahwa integrasi
ketiga komponen tersebut meningkatkan margin keselamatan secara terukur
pada kondisi yang dituju, dengan konsekuensi efisiensi yang dapat
dikuantifikasi dan dipertimbangkan, bukan konsekuensi yang tidak terkendali.
Penelitian ini tidak mendukung klaim bahwa metode usulan berhasil sempurna
atau selalu lebih baik pada seluruh kondisi, mengingat peningkatan yang
tidak merata dan satu kombinasi metrik yang tidak menunjukkan perbedaan
signifikan.

\subsection{Menjawab Pertanyaan Penelitian}

Empat pertanyaan yang telah dirumuskan pada Subbab~\ref{sec:interpretasi}
dapat dijawab sebagai berikut. Pertanyaan mengenai peningkatan clearance
pada skenario dengan kucing pada \textit{blind-spot}, sebagai validasi
hipotesis utama penelitian, didukung kuat oleh data. Pertanyaan mengenai
\textit{trade-off} waktu tempuh pada skenario kompleks juga terjawab,
\textit{trade-off} tersebut memang ada dan cukup besar, bukan sekadar
perubahan kecil yang dapat diabaikan. Pertanyaan mengenai kehalusan gerak
melalui kecepatan angular hanya terjawab sebagian, karena metode usulan
tidak secara seragam menghasilkan kecepatan angular yang lebih rendah pada
seluruh skenario. Pertanyaan mengenai ketiadaan penurunan performa pada
kondisi tanpa \textit{obstacle} terjawab untuk aspek keberhasilan dan
keselamatan, yang tidak menurun sama sekali, meski metrik efisiensi pada
kondisi ini tetap menunjukkan perbedaan statistik, sehingga jawaban atas
pertanyaan ini tidak sepenuhnya bersih dari catatan.

\subsection{Posisi Temuan Penelitian}

Secara keseluruhan, temuan penelitian ini menunjukkan bahwa pendekatan yang
lebih konservatif, dalam arti lebih berhati-hati merespons \textit{obstacle}
dan memperlambat gerak saat diperlukan, dapat meningkatkan keselamatan
navigasi secara terukur. Konsekuensinya terhadap efisiensi juga nyata dan
tidak dapat diabaikan begitu saja. Posisi ini menempatkan metode usulan
bukan sebagai solusi yang unggul pada seluruh aspek, melainkan sebagai
satu titik keseimbangan tertentu antara keselamatan dan efisiensi, yang
kelayakannya pada akhirnya bergantung pada kebutuhan aplikasi yang dituju,
apakah keselamatan ekstra tersebut sepadan dengan tambahan waktu tempuh
yang menyertainya.

\subsection{Keterbatasan Ringkas}

Sebagai pengingat sebelum menutup BAB IV, hasil-hasil di atas diperoleh
sepenuhnya dari simulasi Gazebo dan belum diuji pada robot fisik, serta
belum melalui \textit{ablation study} yang dapat memisahkan kontribusi
masing-masing komponen metode usulan secara independen. Pembahasan lebih
lengkap mengenai keterbatasan ini telah disajikan pada
Subbab~\ref{sec:pembahasan_terintegrasi}.

Hasil eksperimen dan pembahasan pada BAB IV ini telah memberikan gambaran
menyeluruh mengenai performa metode usulan, mencakup keberhasilan,
keselamatan, efisiensi, dan kualitas gerak, beserta konteks dan batasannya.
BAB selanjutnya menyajikan kesimpulan akhir penelitian serta saran untuk
pengembangan lebih lanjut, berdasarkan seluruh temuan yang telah dirangkum
pada subbab ini.